\section{The physical fibre and its ordered cover}
\label{rw:physical-fibre}

Fix an original field and follow its returning arcs outwards. An outer
closing orbit bounds the line-breaking coefficient; three equilibrium-free
strips then supply all intermediate returns, including the nonclosing
ones. We order their local descriptions and transport this one coefficient
inequality to the common interval used in Section~12.
Figure~\ref{rw:physical:interval-scheme} shows the regular and terminal
parts. Their open overlap carries the same $K$ and permits a seam beyond
a preceding regular zero.

\begin{figure}[htbp]
\centering
\begin{tikzpicture}[x=1.25cm,y=.85cm,>=Stealth,font=\small]
 \draw[->] (0,0)--(8.3,0) node[right] {outward input};
 \draw[thick] (0,.75)--(4.6,.75);
 \draw (0,.60)--(0,.90) (4.6,.60)--(4.6,.90);
 \node[above] at (2.3,.75) {$R$: regular initial chain};
 \draw[thick] (3.8,-.75)--(8,-.75);
 \draw (3.8,-.90)--(3.8,-.60) (8,-.90)--(8,-.60);
 \node[below] at (6,-.75) {$T$: terminal union, when present};
 \draw[dotted] (3.8,-1.0)--(3.8,1.05);
 \draw[dotted] (4.6,-1.0)--(4.6,1.05);
 \node[fill=white,inner sep=1pt] at (4.2,0) {$O$};
 \draw[thick,dashed] (3.8,-1.55)--(5.25,-1.55);
 \node[below,align=center] at (4.525,-1.55)
    {possible constant-sign\\connection inside $T$};
\end{tikzpicture}
\caption{Schematic intervals on one actual physical fibre with an outer
closing witness. $O$ is an open overlap. Near the parabolic corner $R$
includes the finite connection, wedge and regular secondary part.
The depicted terminal lower-loss connection, when present, remains
inside $T$. The lengths
carry no metric information, and the arrow orders inputs rather than
trajectories. Proposition~\ref{rw:ordered-cover} supplies these domains.}
\label{rw:physical:interval-scheme}
\end{figure}

We use the original and normalized generators
\[
 \tau_0=a+c,\quad \sigma_0=c(B+m),\qquad
 \tau_1=\bar a+\bar c,\quad
 \sigma_1=\bar c(\bar B+\bar m),
\]
so $\tau_0=\rho\tau_1$, $\sigma_0=\rho^3\sigma_1$,
$\epsilon=\rho^2$ and $\ell=\rho\nu$. Except when explicitly
discussing $\rho=0$, all normalized coordinates have $\rho>0$.
Reflect the whole field once, if necessary, to arrange $\ell\ge0$.
The new outgoing endpoint is the old incoming endpoint; reflection
does not preserve a previously selected numerical outgoing cutoff.
It preserves the original generator norms and cycle multiplicities
by Proposition~\ref{rw:geometry:reflection}.

\subsection{Cycle candidates, retained scales, and outer forcing}

Write
\[
 F_\pm(X)=1\pm\bar aX+\bar BX^2,\qquad
 \mathcal P=\{\bar c=\bar m=0,\ \bar B>0,\
           |\bar a|\le2\sqrt{\bar B}\}
\]
on the normalized max-norm face. At $\mathcal P$ one has
$\bar B\ge1/4$. Indeed either $\bar B=1$ or $|\bar a|=1$,
and the latter inequality gives $\bar B\ge1/4$.
The compactified interval $[0,\infty]$ is always given its ordinary
topology through $p\mapsto p/(1+p)$.
We first cover closing orbits, so that an outer orbit can supply a
coefficient inequality. We will extend the cover to the intervening
comparison arcs after proving that all of those arcs return.

\begin{lemma}[Refined candidate space]\label{rw:fibre:candidate-space}
The limiting data of parameter-dominant counted cycles have a compact
refinement with the following properties.
Away from $\mathcal P$ its coordinates are the normalized
coefficients and the two compactified actual endpoints.
Near $\mathcal P$ it additionally retains the secondary direction
$(C,M)$, the secondary scale and actual first-exit tags.
Every point has one of the local descriptions of Sections~7--10,
including their boundary faces. The same compactification can be
applied simultaneously to an outer cycle and any returning upper
arc between fixed positive initial endpoints and its endpoints.
The assertion about such additional arcs is compactness only;
their existence and local classification will be proved below.
\end{lemma}
\begin{proof}
Let $\chi>0$ be the fixed strong-entry constant in
Lemma~\ref{rw:finite:action-order}. For parameter-dominant
cycles, $p\ge\chi/2$, $q\ge\chi$. Take the intersection, as
$r\downarrow0$, of the closures of their data with $0<\rho<r$
in
\[
 \mathcal S\times[\chi/2,\infty]\times[\chi,\infty],
 \qquad
 \mathcal S=\{(\bar B,\bar m,\bar a,\bar c):
          \max(|\bar B|,|\bar m|,|\bar a|,|\bar c|)=1\}.
\]
This is a compact set. Its points satisfy the nonnegative lower
conditions of Lemma~\ref{rw:finite:lower-cones}. The parameter
$\nu$ tends to zero by Proposition~\ref{rw:geometry:flatness}.

Only the corner at $\mathcal P$ needs further coordinates. Put
\[
 \bar d=\bar c+\nu,\qquad
 \lambda=\max(|\bar d|,\sqrt{|\bar m|}),\qquad
 C=\bar d/\lambda,\quad M=\bar m/\lambda^2
 \quad(\lambda>0).
\]
The direction belongs to the compact face
$\mathcal D=\{\max(|C|,\sqrt{|M|})=1\}$. Use the fixed
sections of Lemma~\ref{rw:secondary:routing}. Retain a finite
flag specifying a wedge exit, a secondary exit, or their overlap.
For a wedge intersection retain $H\in[0,H_0]$ and, if
$H+\lambda>0$, the ratio
$\chi_{\rm exit}=H/(H+\lambda)$.
For a secondary intersection retain its actual radial coordinates
$R_{\rm in},R_{\rm out}\in[0,R_{\max}]$ on the sections
$\phi=\pm\lambda L$, together with
\[
 t=1/P,\qquad \widetilde t=1/Q,\qquad
 0\le t,\widetilde t\le1/p_* .
\]
Here $P,Q$ are the auxiliary linear-flow labels of
Lemma~\ref{rw:secondary:tags}, evaluated at those actual section
points; in particular $t=R/\bar x(T)$ on the incoming section.
They are not defined to be $\lambda p,\lambda q$.
Use a fixed placeholder for an inactive coordinate. Allow all the
indicated flags on an overlap, and take the closure of the resulting
actual-data graph in the finite union of these compact products,
also retaining $\lambda\in[0,\lambda_*]$.
Taking the closure retains radial limits without inverting a map
at zero radius. An orbit in an overlap may have two sets of tags;
keeping both copies preserves compactness and suffices for a finite
subcover.

If $H/\lambda$ tends
to infinity, routing is in the wedge and
$|\bar d|/H+|\bar m|/H^2$ tends to zero. For a finite ratio
the first-exit flag is retained: a first wedge exit has
$H/\lambda\ge L$ and hence the fixed wedge inequalities,
whereas a first secondary exit is in the secondary rectangle.
The equality face has the open overlap proved below. Radial escape is
$R=0$, equivalently $1/P=0$ or $1/Q=0$, and is retained in
the secondary infinity charts. The cases of a finite upper
threshold and its repeated limit are retained with their actual
deficits. Proposition~\ref{rw:secondary:exhaustion} classifies the
entire normalized face, including its repeated and weak boundaries,
and excludes another zero corner. For $\lambda=0$ an actual field
uses only the wedge; secondary directions there record limits with
$\lambda>0$. This classification applies to returning upper arcs.

For completeness the primary alternatives follow directly from
the upper linear flow. A finite nonregular pair is $(p_c,q_c)$
with $\bar m>0$. One finite and one infinite endpoint is a
strict real threshold, a repeated threshold with its whole interval
$p\ge P$, or the weak boundary $\bar m=0$; reflection gives
the other orientation. At two infinite endpoints the strict barrier
of Lemma~\ref{rw:finite:strict-barrier} gives
$\bar c^2+4\bar m\le0$. Strict focus, repeated upper root and
the cubic horizontal boundary are covered in Section~8; the joint
repeated center is Section~9. If $\bar m=0$, then $\bar c=0$
and the lower cones put the point in $\mathcal P$.
At a finite lower loss the first root, its double limit, or the
simultaneous pair gives the alternatives of Section~7.
The two-root and one-infinite statements include roots that escape
to infinity, so these faces are not omitted.

For an intermediate arc, take the product with a second set of actual
tags, sharing the coefficients but retaining its own section points
and first-exit flag.
Use independent initial bounds $p\ge p_0>0$, $q\ge q_0>0$ for
this copy; $q_0$ is uniform on the compact regular entrance.
Impose $p\le p^{\rm out}$, $q\le q^{\rm out}$ before
closing the actual-data graph. This remains a closed subset of a
compact product. The independent tags are needed because an
intermediate arc may take a wedge exit although the outer cycle takes
a secondary exit.
\end{proof}

\begin{lemma}[Uniform selection over compact fibres]
\label{rw:fibre:compact-selection}
In a compact metric product $X$ with continuous outer projection
$\pi:X\to B$ to a compact metric space, let
$X_r$ be the actual data with $0<\rho<r$ and put
$X_0=\bigcap_{r>0}\overline{X_r}$. If a finite union
$U$ of prepared open pieces contains the limiting fibre
$X_0\cap\pi^{-1}(\{b\})$, then
an outer neighborhood and a positive original scale cutoff can be
chosen so that every actual datum over that neighborhood belongs to
$U$. If this construction applies at every outer limit, finitely many
such choices give common cutoffs and bounds.
\end{lemma}
\begin{proof}
Choose $V$ with $\overline V$ disjoint from the compact set
$\pi(X_0\setminus U)$. If no scale cutoff worked, excluded actual data
over $V$ at scales tending to zero would have a subsequential limit in
$(X_0\setminus U)\cap\pi^{-1}(\overline V)$, a contradiction.
A finite outer subcover gives the common minimum scale cutoff and
maximum bounds. The same argument applies to the one-copy graph.
\end{proof}

\begin{lemma}[An endpoint cutoff chosen after contraction]
\label{rw:fibre:candidate-cover}
There is $p_\sharp>0$, chosen before the primary scale cutoff,
and finitely many local candidate neighborhoods that contain every
counted cycle with $p\ge p_\sharp/16$ at sufficiently small
original parameter. The regular member contains all sufficiently
small endpoint data with
$p_\sharp/16\le p\le p_\sharp/2$.
Each other member is either separated from its local center set
on its closure or has a common center-comparison domain and a
positive contraction exponent $b_j>p_\sharp^2$.
Secondary center estimates retain any prescribed finite power
of $\lambda$, or of the wedge height, before division.
\end{lemma}
\begin{proof}
Start with the nonregular part of the compact refined limit set
in Lemma~\ref{rw:fibre:candidate-space}. At a noncenter point
shrink the closed coefficient neighborhood away from the center
set. At a center point the classification gives a divided regular,
two-root, finite-saddle, weak, focus or joint chart.
There are no additional unprepared center losses: the exact
quasi-invariant \eqref{rw:finite:quasi-invariant} implies that an
incoming finite root on a center also has an outgoing root; a
one-infinite center has
$\bar a=-\bar c,\bar m=-\bar B$, so its finite lower factor
equals the upper entry factor; and a lower double root in the
secondary corner has $|\bar a|=2\sqrt{\bar B}>0$ and is off
the secondary center. The repeated coincidence is precisely the
joint chart.

On each center neighborhood, first retain a fixed lower subsegment
with positive action. After reducing the coefficient neighborhood,
normal contraction on this subsegment bounds the absolute coefficient
derivatives of the complete central deviations by
$C\epsilon^{-N}e^{-b_j/\epsilon}$, with $b_j>0$.
These are estimates of the absolute deviations; the logarithm of an
input-independent slow contraction may diverge. The large-power
preparation extends the same estimate by zero up to a lost exit.
Section~5 supplies the real center and the whole input domain on which
to divide. Double division uses two ordinary coefficient derivatives
of the absolute difference, and a single component uses one. The
budgets retained in Sections~7--9 include these derivatives; this step
takes no further input derivative.
In the corner use the weighted absolute estimates of
Lemma~\ref{rw:secondary:flat-budget}. For the present argument
retain power two before the coefficient division; if a later local
coordinate requires more, retain its specified finite power first.
The exact exponent in Proposition~\ref{rw:exact-long} makes
these choices uniform without bounding $\epsilon|\log\lambda|$.
For the wedge choose its fixed widths $|\gamma|,|J|<\delta$
inside the regular-lower center box whenever $|\bar a|$ is
small. The remaining wedge data have $|\bar a|$ bounded
below and are primary off-center data. Thus an off-center
secondary direction attached to a first wedge exit is not
used to estimate that exit by a power of $\lambda$; its
actual radial power is a power of $H$.

Take a finite cover of this refined nonregular set, and then choose
$p_\sharp\le\chi/2$ so that $p_\sharp^2<\min_j b_j$,
and decrease it to meet the small-endpoint error bounds in the
source and regular-column preparations. These choices precede
the smallness choice for $\epsilon$.
The compact complement of the chosen open neighborhoods is
strictly regular. Enlarge that compact regular set to contain
$p_\sharp/16\le p\le p_\sharp/2$ and the limits of cycles
in the two initial energy classes with $p\ge p_\sharp/16$.
The latter enlargement is regular: the source endpoints satisfy
$p=\sqrt{2\epsilon h e_+}\le2\sqrt\eta$ and
$q/p=\sqrt{e_-/e_+}=1+O(\eta)$, while the lower and upper
thresholds have a fixed positive lower bound. The compact-energy
class has $p,q\to0$ and causes no further endpoint limit.
The regular upper return and the positive lower polynomials
therefore give a common compact tube for this enlargement.

Apply Lemma~\ref{rw:fibre:compact-selection} to this one-copy
refined graph. Positive endpoint limits from the initial classes lie
in the regular enlargement, with $q$ bounded below when
$p\ge p_\sharp/16$ by the same source estimates. Thus the selected
finite family covers all the asserted cycle data.
\end{proof}

The candidate cover locates each outer closing orbit. To obtain one
sharp count, we next extract its coefficient inequality; saturation
then supplies all intervening comparison paths.

\begin{proposition}[Forcing supplied by an outer closing zero]
\label{rw:outer-forcing}
Choose the field orientation with $\ell\ge0$.
After a common original parameter shrink, a counted cycle with
$p\ge p_\sharp$ either lies in a constant-sign local comparison
whose full remaining physical passing interval is covered by that
same comparison, or supplies the coefficient inequality
\begin{equation}\label{rw:fibre:original-forcing}
 |\ell|\le e^{-p_\sharp^2/(20\epsilon)}
                    (|\tau_0|+|\sigma_0|).
\end{equation}
Equivalently,
\begin{equation}\label{rw:fibre:normalized-forcing}
 |\nu|\le e^{-p_\sharp^2/(20\epsilon)}
                    (|\tau_1|+\epsilon|\sigma_1|).
\end{equation}
An interior constant-sign piece is not sufficient for the first
alternative. On a center fibre with a closing zero there are no
isolated closing zeros.
\end{proposition}
\begin{proof}
Apply the candidate cover to the outer cycle. A constant-sign chart
containing the entire remaining passing interval gives the first
alternative. In every other chart we derive the coefficient inequality,
including on constant-sign pieces separated from the center set.

The nonnegative lower cone and
\eqref{rw:finite:action-budget}, applied at $p_\sharp$, give
\[
 s\le \exp[-p_\sharp^2/(5\epsilon)].
\]
Indeed
$A_+(p_\sharp)\ge
p_\sharp^2/[2(1+p_\sharp+p_\sharp^2)]$ and the error in
the budget is absorbed after $p_\sharp$ is fixed.
Since $|\ell|<3s$,
any closed primary neighborhood on which
$|\tau_1|+|\sigma_1|\ge g_0>0$ satisfies
\eqref{rw:fibre:original-forcing}: use
$|\tau_0|+|\sigma_0|\ge\rho^3g_0$ and absorb this fixed
power of $\rho$ in the strict gap between the two exponents.

At an off-center secondary first exit, the primary action estimate
alone does not retain the vanishing secondary weight. We use the exact
long match to recover it.
Suppose first $|C|\ge g_0/2$, where
$|\bar a|+|C|\ge g_0$ on the chosen closed neighborhood.
The exact long match gives an outgoing horizontal input
$u_+\le C_1\lambda^{k/b}$, with
$b=1-\epsilon\bar B$, $k=1-2\epsilon\bar B$.
The lower common power has exponent $b/(\epsilon\bar B)$.
The bounded scaled action and scaled prefactor, including a
possible additional contracting lower double passage, imply
\begin{equation}\label{rw:fibre:weighted-budget}
 s\le I_{\rm out}
       \le e^{C_0/\epsilon}\lambda^{k/(\epsilon\bar B)}.
\end{equation}
The first inequality uses the actual reference
$s_+=-\ell H_+\le0$. Choose $\kappa_*>0$ with
$k/\bar B\ge\kappa_*$, then $\lambda_0>0$ so that
\[
 (\kappa_*/2)\log(1/\lambda_0)
                    \ge C_0+p_\sharp^2/5,
 \qquad \epsilon\le\kappa_*/4.
\]
For $\lambda\le\lambda_0$,
\eqref{rw:fibre:weighted-budget} yields
$s\le\lambda^2e^{-p_\sharp^2/(5\epsilon)}$.
For $\lambda\ge\lambda_0$ the primary budget gives the same
bound with the fixed factor $\lambda_0^{-2}$. Consequently
$|\nu|\le\lambda|C|/2$ at small scale and
\[
 |\bar c|=|\lambda C-\nu|\ge\lambda g_0/4,\qquad
 |\sigma_0|\ge c\rho^3\lambda g_0 .
\]
The extra factor $\lambda^2$ proves
\eqref{rw:fibre:original-forcing} uniformly, even if
$\epsilon|\log\lambda|$ is arbitrarily large.
If $|C|<g_0/2$, then $|\bar a|\ge g_0/2$; the corner
has $|\bar c|<g_0/4$, so this reduces to the primary
off-center estimate.

On a retained center neighborhood, cancellation on the real center
recovers the generator factor. Write the complete closing function as
$Z-\rho\nu H$, with $H>h_*$ independent of the section input.
The actual center division gives
\[
 |Z|\le C\epsilon^{-N}e^{-b_j/\epsilon}
             (|\tau_1|+|\sigma_1|+|\nu|).
\]
The local inverse charts and the two ideal comparisons in
Proposition~\ref{rw:generator-division} justify the indicated
generators. In secondary variables retain the weighted version
before returning to primary coefficients: a term
$\lambda^2|C|$ is bounded by
$|\bar d|=\lambda|C|$, and the wedge version is identical
with $H^2|\gamma|\le|\bar d|$ and its retained height.
The small-$|\bar a|$ wedge box was included in these center
comparisons independently of its secondary direction. Since
$|\bar d|\le|\bar c|+|\nu|$ and
$\bar B+\bar m>1/8$ here, the resulting bound is again the
displayed primary one. No inverse power of $\lambda$ is
introduced.
At a closing zero Lemma~\ref{rw:centers:relative-forcing}
absorbs the $\nu$ term. Since $b_j>p_\sharp^2$, the fixed
powers of $\epsilon$ needed to return to original generators
are absorbed with room to spare.

Finally consider the compact regular member at $p\ge p_\sharp$.
Retain a slightly larger actual tube with $p\ge p_\sharp/2$.
The regular upper map has $q\ge p(1-p)$ at small endpoints;
hence both retained actions have a lower bound greater than
$p_\sharp^2/10$ after $p_\sharp$ is reduced. The regular
absolute maps and double center division yield
\[
 |\ell|\le C\epsilon^{-N}e^{-p_\sharp^2/(10\epsilon)}
                      (|\tau_0|+|\sigma_0|),
\]
which implies the required exponent $p_\sharp^2/20$.
The finite candidate cover makes these choices uniform. If the
generators vanish, absorption first gives $\nu=0$, and the actual
center identity excludes isolated zeros.
\end{proof}

\subsection{Equilibria and the three returning strips}

To continue the comparison paths between an initial input and an outer
cycle, we must exclude an equilibrium that could trap an intervening
orbit. We do this first, then fill the upper and two lower strips by
trajectories.

\begin{lemma}[Equilibria enclosed by counted cycles]
\label{rw:fibre:enclosed-equilibria}
On a common small original parameter neighborhood, the Jordan disk
of every counted cycle contains the origin and no other equilibrium.
The section points of any two counted cycles belong to one interval
of the actual return domain.
\end{lemma}
\begin{proof}
The origin has index one and remains the only equilibrium in a
fixed inner disk. The fixed regular collar strips remain free of
equilibria. Proposition~\ref{rw:collar} places every counted
cycle essentially in the collar, surrounding that disk. The upper
vertex calculation in that proposition leaves at most one additional
equilibrium in the upper chart, and it is a saddle.
It remains to exclude horizontal equilibria from the cycle's disk.

If $\rho=0$, then $B=m=a=c=0$ and $\dot x=-y$; the
equilibrium equations give only the origin. Suppose $\rho>0$.
In a horizontal chart $x=\sigma/r$, $y=v/r$, $\sigma=\pm1$,
the equation $\dot y=0$ is
\[
 \sigma(r+v)+dv^2+crv=0.
\]
Its unique small solution is $v=rh(r)$ with $h=-1+o(1)$.
Thus $|z|=|1+y|<\delta$ at any horizontal equilibrium.
The other equilibrium equation reads
\[
 1-z+\bar BX^2+\bar aX+\epsilon\bar m(1-z)^2=0.
\]
Choose fixed small $R,\delta>0$ and then small $\epsilon$.
The left side is positive for $|X|\le R$, $|z|\le\delta$,
so the equilibrium lies in one of the two half-strips $|X|>R$.

In the incoming half-strip put $T=1/(-X)$ and use the positive
time change $\dd\tau/\dd\xi=-X$. With
$\eta_\ell=\rho\ell$, the exact equations are
\begin{equation}\label{rw:fibre:horizontal-equilibrium-field}
\begin{split}
 T_\tau&=\epsilon T\bigl[\bar B-\bar aT+
        \{1-z+\epsilon\bar m(1-z)^2\}T^2\bigr],\\
 z_\tau&=-z+\epsilon\bar cTz(z-1)
                             +\eta_\ell T(1-z)^2 .
\end{split}
\end{equation}
On $0\le T\le1/R$, $|z|\le\delta$, the tangential derivatives
are $O(\epsilon)$ and the normal derivative is
$-1+O(\epsilon+|\eta_\ell|)$.
Extend to a slightly larger rectangle, preserving $T=0$ and
these strict margins. The normally contracting graph construction
used in Section~4 gives
\[
 z=g(T)=\eta_\ell T\,O_{C^2}(1),\qquad
 Z=z-g(T),\qquad Z_\tau=-\beta(T,Z)Z,\quad\beta>1/2.
\]
The two factors in $g$ follow by solving on $\eta_\ell=0$
and $T=0$, respectively, and dividing these vanishing factors
in the bounded graph family. Every equilibrium lies on $Z=0$.
The small original quantity is $\rho\ell$; no bound on
$\nu=\ell/\rho$ is used.

A counted cycle has $z>0$ by
Proposition~\ref{rw:geometry:cycle-height}. At $T=1/R$,
$0\le z\le\delta$, one has $T_\tau>0$, while
$|g|<\delta/2$. A cycle point with $Z<0$ cannot leave the
rectangle in backward time through $T=1/R$; it cannot reach
$z=\delta$ because $Z<0$ is preserved and $g<\delta/2$.
It cannot reach $z=0$, and its compact affine orbit stays away
from $T=0$. It would therefore remain in the rectangle for
all backward time, whereas $|Z|$ grows exponentially.
This is impossible. A periodic orbit on $Z=0$ would be a
nonconstant periodic orbit of the scalar equation on that graph,
also impossible. Every cycle point in this half-strip consequently
lies strictly above $g$.

An equilibrium with $0<z=g(T)<\delta$ is joined vertically to
$z=0$ without meeting the cycle: cycle points with that $T$
are above $g$, including any points above the rectangle.
The line $z=0$ is outside the cycle's disk, since the whole
cycle has $z>0$. Thus the equilibrium is outside the disk.
The same holds immediately if its height is nonpositive.
Reflecting the field and reversing time gives the argument on
the outgoing half-strip, for either sign of $\ell$.
The coordinate $T$ includes finite roots and roots escaping
to infinity, so no horizontal escape case remains.

The single upper saddle is now the only possible additional
equilibrium inside the disk. Its index would change the enclosed total
from one to zero, contradicting the index of a periodic orbit. Thus
only the origin is enclosed. Two counted cycles are nested essential
curves and bound an equilibrium-free annulus in the collar. Cut this
annulus at the transverse central section. The compact-limit argument
for the strips below shows that every intermediate section point
returns, placing both cycles in one interval of the actual return
domain.
\end{proof}

\begin{proposition}[Saturation of the physical interval]
\label{rw:saturation}
Assume \eqref{rw:fibre:original-forcing}, fix
$p_0\in[p_\sharp/8,p_\sharp/4]$, and let $\gamma$ be a counted
cycle with $p_\gamma>p_0$. Every input
$p\in[p_0,p_\gamma]$ has an actual upper return $q(p)$
and both actual lower central intersections. They satisfy
\[
 q'(p)>0,\qquad L_+'(p)<0,\qquad L_-'(q)<0,
 \qquad
 L_+(p),\ L_-(q(p))\ge s_\gamma>0.
\]
All three arcs, and the segment joining their central endpoints,
lie in the fixed collar. This statement holds at nonclosing inputs
as well as on periodic orbits.
\end{proposition}
\begin{proof}
At $p_0$ the fixed small-endpoint upper map and both regular
lower maps exist. After reducing $p_\sharp$, their two actions
are at most $p_\sharp^2/31$ uniformly in the stated interval
for $p_0$. Their positive homogeneous amplitudes are bounded
below by a fixed positive prefactor times
$e^{-p_\sharp^2/(31\epsilon)}$.
The actual signed reference contributions are $O(|\ell|)$.
Since the original generator norm is bounded and the forcing
exponent is $p_\sharp^2/20$, both complete central heights are
positive for sufficiently small scale. Their arcs lie in the collar
by the regular construction.

The axis endpoints of the three initial arcs lie inside the
corresponding endpoints of $\gamma$. Since trajectories cannot cross
$\gamma$, the arcs stay in its disk. Join their central endpoints by
the central crosscut. Each lower path stays in its own half-plane until
its first central intersection, so it meets the crosscut only there.
The three arcs and the crosscut form a simple essential curve in the
collar, inside $\gamma$, whether or not its two central values coincide.

The initial and outer upper arcs, joined by the two axis segments,
bound a compact equilibrium-free strip by
Lemma~\ref{rw:fibre:enclosed-equilibria}. Since $\dot y=x$ on
$y=0$, the positive axis segment is an entrance, the negative one
an exit, and the orbit sides cannot be crossed. A nonexiting orbit
would have an equilibrium-free compact $\omega$-limit set, hence a
periodic orbit by Poincar\'e--Bendixson. Its Jordan disk lies in
$y>0$ inside $\gamma$ and contains no equilibrium, contradicting
the periodic-orbit index. Thus every intermediate upper entry returns, with
$q_0\le q(p)\le q_\gamma$. The transverse flow determinant
has a fixed sign and gives $q'(p)>0$.

Construct the right and left lower strips from the corresponding
initial and outer arcs, an axis segment, and a central crosscut.
They are equilibrium-free as well. On the right use backward
time, so the positive axis segment is an entrance; on the left
use forward time and the negative axis segment is an entrance.
The central crosscuts are exits because $\dot x>3/4$ on the
fixed central rectangle. The same compact-limit and index argument
gives both central intersections at every intermediate endpoint.
Transversality and endpoint order give $L_+'(p)<0$, $L_-'(q)<0$;
their minimum heights occur on $\gamma$ and equal $s_\gamma>0$.

The three strips and the joining central segments lie in the annulus
between the two essential curves, hence remain in the fixed collar.
The nonexit contradiction uses compactness for each fixed field;
flight times may diverge as a saddle or slow gap tends to zero.
\end{proof}

On this physical interval keep the actual central height as output.
The fixed weak references give
\[
 F(p)=L_+(p)-L_-(q(p))=U(p)-V(p)-\rho\nu H,\qquad H>0,
\]
where $H$ is independent of $p$ and $U',V'<0$. Consequently
\begin{equation}\label{rw:fibre:physical-log}
 K=\epsilon\log(V'/U'),\qquad
 F'=U'(1-e^{K/\epsilon}).
\end{equation}
A common regular input change and a change of the fixed additive
references preserve $K$ on this orbit interval. We keep the original
central height as output, since a nonlinear output change need not
preserve it. The signed center-comparison domains used to estimate
these functions are auxiliary; the physical interval of
Proposition~\ref{rw:saturation} has positive complete central heights.

\subsection{The ordered cover of all intermediate arcs}

Saturation supplies the intermediate paths. We must now order their
local descriptions so that separate endpoint bounds become a bound
on one interval. The whole-passage tests will give a regular prefix
followed by at most one terminal union. At the parabolic corner this
requires two connections: from the finite family to the wedge, and
from the wedge to the secondary family.

\begin{lemma}[The two parabolic connections]
\label{rw:fibre:parabolic-connections}
Near $\mathcal P$ the finite regular family connects by an actual
open overlap to the wedge, possibly through a constant-sign lower
loss interval. The wedge in turn has an actual open overlap with
the regular secondary family. The construction applies to the
nonperiodic returning arcs in Proposition~\ref{rw:saturation}.
\end{lemma}
\begin{proof}
Consider first the incoming lower double boundary
$\bar a_*=2\sqrt{\bar B_*}$ and put
$r=1/\sqrt{\bar B_*}$. The limiting upper return has $q=p$.
For $0<p<r$,
\[
 A_-(p)-A_+(p)
 =2\bar a_*\int_0^p
               \frac{X^2}{F_+(X)F_-(X)}\,\dd X>0.
\]
Thus $K<0$ on every fixed compact regular initial interval,
uniformly in a smaller off-center neighborhood.
Let $C_w=\sqrt{1+2w_0}/w_0$. The limiting upper parabola
$Y=(p^2-X^2)/2$ meets $w=w_0$ at $H=C_w/p$.
Choose $L_1>2r$ so large that $C_w/L_1<H_0/2$.
The upper arcs for $L_1\le p\le2L_1$ form a compact regular
family; their transverse intersections persist and belong to
the wedge after the coefficients are reduced.

The incoming action tends to infinity as its first double root
is approached from below. Choose $R<r$ with
\[
 A_-(R)>\max_{p\le2L_1}A_+(p)+3,
\]
where the maximum is on the retained regular outgoing interval.
Then choose a fixed $p_1<r$ with $p_1>R+2\zeta$ for some
$\zeta>0$. Uniform convergence of the regular upper map to
$q=p$ on $[p_1,2L_1]$ gives $q(p)>R+\zeta$ there.
Lemma~\ref{rw:finite:incoming-contraction} now gives $K<0$
on its entire physical passing part, even beyond the lower
root where the ordinary regular action formula no longer applies.
This interval overlaps both the regular initial family and the
wedge. The corresponding lower slow action in the wedge has
the same strict sign and dominates its bounded remainder.
Reflection gives the other lower double boundary and reverses
the sign. Away from these boundary neighborhoods the lower
factors have a common positive margin on every retained finite
interval, so the regular family connects directly to the wedge.

For the wedge--secondary overlap, choose the fixed $L$ in secondary
routing large enough before decreasing any scale. Require
$2/L<\delta$ and $4/L^2<\delta$. At
$w=w_0$, $H=k_0\lambda L$, with $1/2<k_0<1$, the exact
secondary coordinates are
\[
 \widehat x=\frac1{k_0Lw_0},\qquad
 \widehat z=\frac1{k_0^2L^2w_0}.
\]
The auxiliary linear flow defining the actual tags gives
$P,Q=O(C_w/L)$ uniformly for bounded $(C,M)$.
At $C=M=0$ the values are $C_w/(k_0L)$; the perturbation
on these short arcs is of higher order in $1/L$.
Increase $L$ so all these tags lie in one compact regular
interval below the finite secondary thresholds.
The wedge inequalities are strict, and the same actual arcs
meet the secondary sections before reaching $w=w_0$.
Transversality persists as $k_0$ varies on a smaller open subinterval
of $(1/2,1)$. Thus an open family of the same returning arcs has both
descriptions, including the nonclosing comparisons.
\end{proof}

\begin{proposition}[An ordered cover on each physical fibre]
\label{rw:ordered-cover}
After a common original parameter shrink, every physical interval
of Proposition~\ref{rw:saturation} is covered by a connected
regular initial part and an ordered union of the local passing
charts of Sections~7--10. Their adjacent members overlap on
actual orbits. At a joint repeated center one contact member and
the full separated member form a single union. Near $\mathcal P$
the initial part includes the finite connection, the wedge and
the regular secondary part.
After grouping overlapping descriptions of the same end, there
is at most one terminal union after this initial part.
It may contain a constant-sign lower-loss connection followed
by a singular upper end; that connection is retained as its
own piece.

The coefficient neighborhoods and all cutoffs can be selected
from a finite family, independently of the chosen outer cycle.
The order includes signed split-root faces and the repeated
one-infinite interval. Each local derivative estimate is applied
on its stated whole-passage domain. This proposition makes no
sum of local zero bounds.
\end{proposition}
\begin{proof}
Fix the finite derivative orders of Sections~7--10 and retain their
larger construction domains before reducing sections or coefficient
neighborhoods. Saturation supplies actual upper and lower returns
with positive central heights on $[p_0,p_\gamma]$.
Since $q'>0$, their endpoints satisfy $p\le p_\gamma$,
$q\le q_\gamma$ and inherit the outer lower cones.

The key ordering fact is that a lost passing component cannot reappear.
The exact lower factor on the fixed phase interval is
\[
 F_I(x)=\bar B-\bar ax+(1+\epsilon\bar m)x^2
 -\epsilon\nu(1+2\epsilon\bar m)x^3H(x)
 +\epsilon^3\nu^2\bar m x^4H(x)^2.
\]
By Proposition~\ref{rw:phase}, it is strictly convex and positive
near the fixed exit $x_R$. Lemma~\ref{rw:transitions:lower-sections}
requires $F_I(s)>0$ at every point between the actual entry and
$x_R$. Thus only the component adjoining $x_R$ passes: if a root
is present, the entry must lie to the right of its last root.
The entry decreases as its actual endpoint increases, since
$T=1/p$ or $1/q$ and $t_T>0$. Its allowed endpoints consequently
form one outward initial interval. Entry positivity in the other
component would not suffice. Also $F_I'(0)=-\bar a$ and its
reflected derivative is $\bar a$; strict convexity permits an
interior positive double minimum on at most one of the two sides.

The primary upper slow tests have the same one-sided form. For a repeated
factor $v^2+\mu_U$, passage to the positive fixed exit requires
$\mu_U>0$ or $v>\sqrt{-\mu_U}$; its actual finite entry decreases
with $p$. This retains the negative entries beyond the repeated
threshold whenever they pass. For the weak factor $u(\gamma+u)$,
the whole domain is $\max(0,-\gamma)<u<u_0$, and
$p=\hat p-t(u)$ with $t'>0$. At joint contact both calibrated
entries decrease, so the two whole-path tests are intersections
of one-sided restrictions. General horizontal factors use the
whole lower test above; a smaller monotone block alone permits
an endpoint test. Signed auxiliary coefficient comparisons do not
remove physical inputs, whose complete central heights are positive
by saturation.

The increasing upper graph and its reflected inverse are those of
Lemma~\ref{rw:finite:upper-return}. Its nonregular ends are classified
in Lemma~\ref{rw:fibre:candidate-space}, with the full repeated
interval $p\ge P$, $q=\infty$ retained. Following the graph outward
gives these terminal templates:
\begin{itemize}
\item At a regular finite upper end, the first lower loss is a
single root or the simultaneous simple-root pair. A simple root
terminates passage; a double retains its signed slow alternatives.
\item At a finite upper saddle, the lower paths are regular, in
one or both contact neighborhoods, or have one incoming double-loss
connection before the saddle. An outgoing double with regular
incoming outer data is excluded by
Lemma~\ref{rw:finite:action-order}, applied to the outer cycle.
The incoming connection is the whole constant-sign interval of
Lemma~\ref{rw:finite:incoming-contraction}.
\item At a one-infinite end, a lower root can precede, coincide
with, or follow the upper threshold. These give respectively the
lower connection, the contact neighborhood, or the upper-tail piece
followed by the lower loss. A root escaping to infinity retains
its finite-root and cubic descriptions in the same family.
\item At two infinite endpoints the lower cones exclude a prior
simple root. Section~8 supplies the strict-focus, one lower-double,
repeated, opposing-slow and cubic alternatives; Section~9 supplies
the joint contact interval and its full separated tail.
\end{itemize}
A finite saddle has one returning quadrant, while the strict,
repeated and weak preparations have the through components just
described. Together with the monotone entry tests, this excludes a
second independent upper end or a return after a split-root
obstruction. The pieces of each listed end therefore form one
terminal union, retaining a preceding constant-sign lower connection
as a separate piece within it.

Near $\mathcal P$, use the actual first exits and tags of
Proposition~\ref{rw:secondary:exhaustion}. Routing confines a
returning arc to the fixed secondary rectangle and uniqueness
excludes another middle excursion. The normalized face
$\max(|C|,\sqrt{|M|})=1$ has the finite-saddle, strict real,
repeated, weak and focus alternatives, including $1/P=0$ and
$1/Q=0$. The repeated row retains its full label interval.
At its weak end the actual tag is $P=P_s-\psi(t(u))$ with
$\psi'>c$; again outward motion decreases $u$.
The preceding lemma inserts the finite connection, wedge and
regular secondary prefix. The same first-threshold reasoning then
gives one secondary terminal family; a lower-double overlay can
add its retained connection but cannot add an opposing double or
another returning quadrant.

Retain actual open seams before reducing the scale: a fixed positive
deficit crosscut between finite regular and singular pieces; the
joint contact/separated overlap $P+d<p<P+2d$ with both gaps
positive; and a fixed reciprocal annulus before either slow block,
using Lemma~\ref{rw:transitions:common-input}. These connectors
are transverse with bounded inverses. The finite lower-double
bridge has the strict action margin of the preceding lemma.
That lemma also supplies the finite--wedge and wedge--secondary
transverse overlaps. Within the secondary family retain its fixed
deficit crosscuts and reciprocal annuli. Adjacent strict, repeated
and weak faces use their own uniform preparations; no estimate at
a fixed positive eigenvalue ratio is extended across zero or infinity.
All seams compare the same original central function.

Apply Lemma~\ref{rw:fibre:compact-selection} to the joint compactification
of Lemma~\ref{rw:fibre:candidate-space}: the intermediate copy has
independent bounds $p\ge p_0$, $q\ge q_0>0$, its own first-exit tags,
and $p\le p^{\rm out}$, $q\le q^{\rm out}$. Each compact limiting
fibre is covered by the ordered pieces and their whole regular prefix.
The retained crosscuts, action margins and transverse seams make this
an open cover, including reciprocal radial ends; the one-sided entry
tests preserve order. Positive endpoint limits from the initial energy
classes use the regular enlargement in
Lemma~\ref{rw:fibre:candidate-cover}. Thus the lemma gives common
cutoffs after all sections, derivative budgets and overlap margins
have been fixed.
\end{proof}

\subsection{Transport in the original coefficient fibre}

We transfer the outer orbit's inequality to the splitting, gap and
saddle-ratio coordinates. The comparisons below may contain $|\nu|$;
the exponential factor absorbs it on the fixed coefficient fibre.

\begin{proposition}[Uniform coefficient and domain transport]
\label{rw:coefficient-transport}
On the ordered cover of Proposition~\ref{rw:ordered-cover},
the fixed original inequality
\eqref{rw:fibre:original-forcing} supplies the relative coefficient
bounds used in the local critical-function estimates, uniformly at
all physical inputs, including
nonclosing inputs. More precisely, for each local center-generator
size $I_{\rm loc}$ used in Sections~7--10 there is a modulus
$\eta_{\rm loc}(\epsilon)\to0$ such that
\[
 |\nu|\le\eta_{\rm loc}(\epsilon)I_{\rm loc}.
\]
The coefficient constants and moduli are common to the finite cover.
The coefficient divisions take place on the larger comparison
domains of Section~5 and the corresponding secondary domains.
The counted inputs remain on the positive-central physical interval.
If a local generator size vanishes, absorption first gives
$\nu=0$ and then the relevant exact center identity.
\end{proposition}
\begin{proof}
Set $\delta_\epsilon=e^{-p_\sharp^2/(20\epsilon)}$.
The original inequality is
\begin{equation}\label{rw:fibre:transport-start}
 |\nu|\le\delta_\epsilon
                 (|\tau_1|+\epsilon|\sigma_1|)
          \le\delta_\epsilon
                 (|\tau_1|+|\sigma_1|).
\end{equation}
We compare the original generators with the local transverse
coefficients in each type of chart.

On regular and strict-focus charts take
$I_{\rm loc}=|\tau_1|+|\sigma_1|$. At a two-lower-root match,
the strict splitting derivative of
Proposition~\ref{rw:finite:two-root-count} makes
$(\bar a,\delta,\nu)$ a bounded inverse chart on a larger box.
Here $\delta$ is the actual strong-intersection splitting.
Single division from its exact reversible axis gives
\[
 |\tau_1|+|\sigma_1|
                 \le C(|\bar a|+|\delta|+|\nu|).
\]
The finite upper saddle uses $(\bar a,\kappa,\nu)$ with the
actual $\kappa=r-1=-\operatorname{tr}/\lambda_+$; the derivative
$\partial_{\bar c}\kappa=-1/\sqrt{\bar m}+O(\epsilon)$ has a
nonzero margin. The same division, already described after
Proposition~\ref{rw:generator-division}, gives
$|\tau_1|+|\sigma_1|\le C(|\bar a|+|\kappa|+|\nu|)$.
Off these center axes the transverse size has a positive lower
bound on each smaller closed chart.

At the finite three-saddle intersection,
\eqref{rw:centers:ideal-comparison} gives
\[
 |\tau_1|+|\sigma_1|+|\nu|
            \asymp |\xi_2|+|\kappa\xi_1|+|\nu|.
\]
Thus take $I_{\rm loc}=|\xi_2|+|\kappa\xi_1|$; the product
retains both intersecting center branches. The double-division
segments and derivative budget are those of Section~5.

At a strict one-infinite simple-root coincidence use the exact
splitting $\delta$ and the actual scaled exponent difference
$g$. On its strict compact row the determinant of this coordinate
change tends to
\[
 \frac{p_c^4}{\bar B(1-\bar Bp_c^2)^3}>0,
\]
up to the positive unit from the regular power coordinate.
The calculation and the equality of the actual leading powers
on the second center are in
Proposition~\ref{rw:infinity:strict-lower-loss}.
Their bounded inverse and single division give
$|\tau_1|+|\sigma_1|\le C(|\delta|+|g|+|\nu|)$.
At the repeated and weak faces use instead the two comparisons
\eqref{rw:centers:single-ideal-comparison}, with local sizes
$|\tau|+|\zeta|$ and $|d_s|+|\varsigma|$, respectively.
Their inverse charts and exact center axes were proved in
Proposition~\ref{rw:center-coordinates}; the repeated comparison
uses \eqref{rw:coordinates:zeta-identity}, rather than the strict-row
determinant. The fixed-mean and lifted-contact comparisons, and the
single weak offset $\bar B+K_*A$ of
Proposition~\ref{rw:domains:weak}, are defined on larger domains
containing every original input and the whole coefficient segments.
At a cubic two-infinite end $\bar B=0,\bar m<0$,
$|\bar B+\bar m|$ has a positive margin. Hence
$|\tau_1|+|\sigma_1|\asymp|\bar a|+|\bar c|$,
the cubic generator size. On other constant-sign pieces
either the same bounded comparison applies or the
transverse size is bounded below; no division at a vanishing
gap is needed.

At the parabolic corner we must retain the scale weights before
comparing generator sizes. Since
$\bar B+\bar m>1/8$, bounded normalized coefficients give
\[
 |\tau_1|+|\sigma_1|
       \le C(|\bar a|+|\bar c|)
       \le C(|\bar a|+|\bar d|+|\nu|).
\]
Inserting this in \eqref{rw:fibre:transport-start} and
absorbing gives
\begin{equation}\label{rw:fibre:weighted-transport}
 |\nu|\le2C\delta_\epsilon(|\bar a|+|\bar d|).
\end{equation}
On the wedge, $\bar d=H\gamma$ with $H\le H_0<1$;
on the secondary charts, $\bar d=\lambda C$ with
$\lambda\le\lambda_*<1$. It follows, uniformly at every
actual input, that the forcing is small relative to the local
sizes $|\bar a|+|\gamma|$ and $|\bar a|+|C|$.
The stronger weighted inequality
\eqref{rw:fibre:weighted-transport} is retained whenever
those local expressions are returned to original coefficients.
The wedge bound in \eqref{rw:secondary:critical-forcing} retains
$|\bar a|+|\bar d|$; the secondary bound uses $|\bar a|+|C|$.
The additional powers in \eqref{rw:secondary:local-forcing} arise
when an actual secondary closing witness supplies its absolute
budget; they are stronger than the bounds used here.

Near a finite secondary upper saddle the exact ratio calculation
\eqref{rw:secondary:saddle-ratio} has a nonzero derivative
in $C$ and yields
\[
 |C|\le C_0\bigl(|\kappa|+
                  \epsilon\lambda|\bar a|
                  +\epsilon\lambda|\nu|\bigr).
\]
Thus the preceding bound implies
$|\nu|\le\widetilde\eta(\epsilon)
 (|\bar a|+|\kappa|)$ after one further absorption.
The estimate therefore uses the actual ratio of the secondary field,
including its coefficient error relative to the limiting linear field.
On the weak and real repeated secondary rows $|C|$ has a
positive margin on the normalized face, so their transverse
size is bounded below. They need no inverse coordinate
division at the vanishing weak eigenvalue. A strict focus
center has its regular division in $(\bar a,C)$, and the
finite saddle center was treated above. These neighborhoods
cover the normalized face with the adjacent strict rows,
so the constants are uniform on their common boundary faces.

In every unweighted comparison we have obtained
$I_{\rm global}\le C(I_{\rm loc}+|\nu|)$.
Choose the common scale with $C\delta_\epsilon<1/2$;
then \eqref{rw:fibre:transport-start} gives
$|\nu|\le2C\delta_\epsilon I_{\rm loc}$.
The weighted comparisons have just given the same conclusion
with their explicitly retained factors. If $I_{\rm loc}=0$,
this inequality gives $\nu=0$ before any division by that
size. The inverse coordinate identity then places the field
on its real center axis, so its actual compared maps agree.
This deals with the zero-size case as well.

These are coefficient inequalities, so they hold also at nonclosing
inputs. All divisions use the larger domains of Section~5, while
saturation supplies positive central heights for the original field.
Taking the largest of the finite coefficient bounds makes the
absorptions and moduli uniform on the whole interval.
\end{proof}

